% GENERATED FILE. Edit canonical lessons, then run npm run generate:books.
% canonical-source-hash: fnv1a64:ac8b10e5891ea007
% canonical-lessons: GU-C501-lakdu, GU-C501-karigar, GU-C501-ojar, GU-C501-avadvu, GU-C501-vyavsay

\chapter{The Craftsman's Bench}
\label{ch:gu-the-craftsman-s-bench}
\hlchaptermodality{501}

\begin{chapteropening}
\textbf{By the end of this chapter:} \emph{I can say what a craftsman makes and what he knows how to do.}

Complete the last lesson of chapter 501: I can say what a craftsman makes and what he knows how to do.
\end{chapteropening}

\section[\texorpdfstring{\gu{લાકડું} (lākḍũ)}{લાકડું (lākḍũ)}]{\gu{લાકડું} (lākḍũ) — wood}

\label{lesson:GU-C501-lakdu}



\begin{quote}
Before the new one: say the Gujarati for to pester, then the Gujarati for to hope.
\end{quote}

\subsection*{You'll want to know: \gu{લાકડું}}
\textbf{\gu{લાકડું}} — \emph{lākḍũ} — \textquotedblleft{}wood\textquotedblright{}.

\begin{quote}
\gu{આ} \gu{ખુરશી} \gu{જૂના} \gu{લાકડાની} \gu{છે}.
\end{quote}

\begin{quote}
\emph{ā khurśī jūnā lākḍānī che.}
\end{quote}

\begin{quote}
This chair is made of old wood.
\end{quote}

A neuter noun.

\begin{cousinweb}
\textbf{\gu{લાકડું}} (\emph{lākḍũ}) goes back to Sanskrit \emph{lakuṭa}, \textquotedblleft{}a club, a stick\textquotedblright{}, and CDIAL lists the Gujarati word with the sense \textquotedblleft{}stick\textquotedblright{} (CDIAL 10875). In everyday Gujarati it names wood itself, the stuff a chair or a door is made of.
\end{cousinweb}

\begin{grammarlens}[title={What you've built}]
One more word for describing what you lived and felt.
\end{grammarlens}

\subsection*{Your turn}
\begin{itemize}
\raggedright
  \item \emph{Say it:} \emph{lākḍũ}
  \item \emph{Say it:} \emph{lākḍũ}, once more
  \item \emph{Say it:} \emph{ā khurśī jūnā lākḍānī che.} — This chair is made of old wood.
  \item \emph{Recall:} say the Gujarati for to pester, then the Gujarati for to hope, then say \emph{lākḍũ} again
\end{itemize}

\begin{tcolorbox}[breakable,colback=teal!4,colframe=teal!35!black,title={Before you move on}]
What does \gu{લાકડું} mean? (\textquotedblleft{}Wood\textquotedblright{}.) Say it once more.
\end{tcolorbox}
\section[\texorpdfstring{\gu{કારીગર} (kārīgar)}{કારીગર (kārīgar)}]{\gu{કારીગર} (kārīgar) — a craftsman}

\label{lesson:GU-C501-karigar}



\begin{quote}
Before the new one: say the Gujarati for to hope, then the Gujarati for wood.
\end{quote}

\subsection*{You'll want to know: \gu{કારીગર}}
\textbf{\gu{કારીગર}} — \emph{kārīgar} — \textquotedblleft{}a craftsman\textquotedblright{}.

\begin{quote}
\gu{આ} \gu{કારીગર} \gu{લાકડામાંથી} \gu{સુંદર} \gu{રમકડાં} \gu{બનાવે} \gu{છે}.
\end{quote}

\begin{quote}
\emph{ā kārīgar lākḍāmā̃thī sundar ramakḍā̃ banāve che.}
\end{quote}

\begin{quote}
This craftsman makes beautiful toys out of wood.
\end{quote}

A masculine or feminine noun, by the person meant.

\begin{cousinweb}
\textbf{\gu{કારીગર}} (\emph{kārīgar}) is usually traced to Persian \emph{kārīgar}, \textquotedblleft{}a craftsman\textquotedblright{}, a word built on \emph{kār}, \textquotedblleft{}work\textquotedblright{}.
\end{cousinweb}

\begin{grammarlens}[title={What you've built}]
One more word for describing what you lived and felt.
\end{grammarlens}

\subsection*{Your turn}
\begin{itemize}
\raggedright
  \item \emph{Say it:} \emph{kārīgar}
  \item \emph{Say it:} \emph{kārīgar}, once more
  \item \emph{Say it:} \emph{ā kārīgar lākḍāmā̃thī sundar ramakḍā̃ banāve che.} — This craftsman makes beautiful toys out of wood.
  \item \emph{Recall:} say the Gujarati for to hope, then the Gujarati for wood, then say \emph{kārīgar} again
\end{itemize}

\begin{tcolorbox}[breakable,colback=teal!4,colframe=teal!35!black,title={Before you move on}]
What does \gu{કારીગર} mean? (\textquotedblleft{}A craftsman\textquotedblright{}.) Say it once more.
\end{tcolorbox}
\section[\texorpdfstring{\gu{ઓજાર} (ojār)}{ઓજાર (ojār)}]{\gu{ઓજાર} (ojār) — a tool of a trade}

\label{lesson:GU-C501-ojar}



\begin{quote}
Before the new one: say the Gujarati for wood, then the Gujarati for a craftsman.
\end{quote}

\subsection*{You'll want to know: \gu{ઓજાર}}
\textbf{\gu{ઓજાર}} — \emph{ojār} — \textquotedblleft{}a tool of a trade\textquotedblright{}.

\begin{quote}
\gu{સુથાર} \gu{પોતાનાં} \gu{ઓજારો} \gu{એક} \gu{પેટીમાં} \gu{રાખે} \gu{છે}.
\end{quote}

\begin{quote}
\emph{suthār potānā̃ ojāro ek peṭīmā̃ rākhe che.}
\end{quote}

\begin{quote}
The carpenter keeps his tools in a box.
\end{quote}

A neuter noun.

\begin{cousinweb}
\textbf{\gu{ઓજાર}} (\emph{ojār}) is usually traced to Arabic \emph{awzār}, a plural meaning \textquotedblleft{}loads\textquotedblright{} or \textquotedblleft{}gear\textquotedblright{}, which reached Gujarati through Persian. The plural is \textbf{\gu{ઓજારો}}.
\end{cousinweb}

\begin{grammarlens}[title={What you've built}]
One more word for describing what you lived and felt.
\end{grammarlens}

\subsection*{Your turn}
\begin{itemize}
\raggedright
  \item \emph{Say it:} \emph{ojār}
  \item \emph{Say it:} \emph{ojār}, once more
  \item \emph{Say it:} \emph{suthār potānā̃ ojāro ek peṭīmā̃ rākhe che.} — The carpenter keeps his tools in a box.
  \item \emph{Recall:} say the Gujarati for wood, then the Gujarati for a craftsman, then say \emph{ojār} again
\end{itemize}

\begin{tcolorbox}[breakable,colback=teal!4,colframe=teal!35!black,title={Before you move on}]
What does \gu{ઓજાર} mean? (\textquotedblleft{}A tool of a trade\textquotedblright{}.) Say it once more.
\end{tcolorbox}
\section[\texorpdfstring{\gu{આવડવું} (āvaḍvũ)}{આવડવું (āvaḍvũ)}]{\gu{આવડવું} (āvaḍvũ) — to know how (to do something)}

\label{lesson:GU-C501-avadvu}



\begin{quote}
Before the new one: say the Gujarati for a craftsman, then the Gujarati for a tool of a trade.
\end{quote}

\subsection*{You'll want to know: \gu{આવડવું}}
\textbf{\gu{આવડવું}} — \emph{āvaḍvũ} — \textquotedblleft{}to know how (to do something)\textquotedblright{}.

\begin{quote}
\gu{મીરાને} \gu{લાકડાનું} \gu{કામ} \gu{આવડે} \gu{છે}.
\end{quote}

\begin{quote}
\emph{Mīrāne lākḍānũ kām āvḍe che.}
\end{quote}

\begin{quote}
Mira knows how to do woodwork.
\end{quote}

Past (it …): \gu{તે} \gu{કામ} \gu{તેને} \gu{આવડી} \gu{ગયું}, \emph{te kām tene āvḍī gayũ}.

\begin{cousinweb}
\textbf{\gu{આવડવું}} (\emph{āvaḍvũ}) takes the person who has the skill with \textbf{\gu{ને}}, and the skill itself is the subject: \textbf{\gu{મને} \gu{તરતાં} \gu{આવડે} \gu{છે}}, I know how to swim. CDIAL lists the Gujarati verb, \textquotedblleft{}to come to someone's knowledge\textquotedblright{}, under Sanskrit \emph{āpatati}, \textquotedblleft{}flies towards; befalls\textquotedblright{} (CDIAL 1193); the step to knowing a skill is not explained there.
\end{cousinweb}

\begin{grammarlens}[title={What you've built}]
One more word for describing what you lived and felt.
\end{grammarlens}

\subsection*{Your turn}
\begin{itemize}
\raggedright
  \item \emph{Say it:} \emph{āvaḍvũ}
  \item \emph{Say it:} \emph{āvaḍvũ}, once more
  \item \emph{Say it:} \emph{Mīrāne lākḍānũ kām āvḍe che.} — Mira knows how to do woodwork.
  \item \emph{Recall:} say the Gujarati for a craftsman, then the Gujarati for a tool of a trade, then say \emph{āvaḍvũ} again
\end{itemize}

\begin{tcolorbox}[breakable,colback=teal!4,colframe=teal!35!black,title={Before you move on}]
What does \gu{આવડવું} mean? (\textquotedblleft{}To know how (to do something)\textquotedblright{}.) Say it once more.
\end{tcolorbox}
\section[\texorpdfstring{\gu{વ્યવસાય} (vyavsāy)}{વ્યવસાય (vyavsāy)}]{\gu{વ્યવસાય} (vyavsāy) — a profession}

\label{lesson:GU-C501-vyavsay}



\begin{quote}
Before the new one: say the Gujarati for a tool of a trade, then the Gujarati for to know how.
\end{quote}

\subsection*{You'll want to know: \gu{વ્યવસાય}}
\textbf{\gu{વ્યવસાય}} — \emph{vyavsāy} — \textquotedblleft{}a profession\textquotedblright{}.

\begin{quote}
\gu{મારા} \gu{કાકાનો} \gu{વ્યવસાય} \gu{સુથારનો} \gu{છે}.
\end{quote}

\begin{quote}
\emph{mārā kākāno vyavsāy suthārno che.}
\end{quote}

\begin{quote}
My uncle is a carpenter by profession.
\end{quote}

A masculine noun.

\begin{cousinweb}
\textbf{\gu{વ્યવસાય}} (\emph{vyavsāy}) is a tatsama, taken from Sanskrit \emph{vyavasāya}, \textquotedblleft{}strenuous effort; resolve\textquotedblright{}, which Monier-Williams also gives as \textquotedblleft{}business, trade\textquotedblright{} (Monier-Williams).
\end{cousinweb}

\begin{grammarlens}[title={What you've built}]
One more word for describing what you lived and felt.
\end{grammarlens}

\subsection*{Your turn}
\begin{itemize}
\raggedright
  \item \emph{Say it:} \emph{vyavsāy}
  \item \emph{Say it:} \emph{vyavsāy}, once more
  \item \emph{Say it:} \emph{mārā kākāno vyavsāy suthārno che.} — My uncle is a carpenter by profession.
  \item \emph{Recall:} say the Gujarati for a tool of a trade, then the Gujarati for to know how, then say \emph{vyavsāy} again
\end{itemize}

\begin{tcolorbox}[breakable,colback=teal!4,colframe=teal!35!black,title={Before you move on}]
What does \gu{વ્યવસાય} mean? (\textquotedblleft{}A profession\textquotedblright{}.) Say it once more.
\end{tcolorbox}
